% GENERATED FILE. Edit canonical lessons, then run npm run generate:books.
% canonical-source-hash: fnv1a64:33badfca4ae65250
% canonical-lessons: FR-C03-merci, FR-C03-de-rien, FR-C03-aller, FR-C03-comment-ca-va, FR-C03-comment-registers, FR-C03-comme-ci-comme-ca, FR-C03-practice

\chapter{How Are You?}
\label{ch:how-are-you}
\hlchaptermodality{3}

\begin{chapteropening}
\textbf{By the end of this chapter:} \emph{I can ask how someone is in French, answer naturally, and return the question.}

Now the small talk that follows a name: thanking, answering a thank-you, and asking how someone is doing. French asks that last question with a verb of motion — not “how are you?” but “how does it go?”
\end{chapteropening}

\section[merci]{merci — \textquotedblleft{}thank you,\textquotedblright{} and it means \emph{mercy}}

\label{lesson:FR-C03-merci}



\begin{quote}
You can introduce yourself in French. Next you'll ask \emph{how are you?} — and every polite answer needs \textbf{merci} first. It hides a word you already know: \emph{mercy}.
\end{quote}

\begin{sounds}
\begin{itemize}
\raggedright
  \item \texttt{r-uvular} — \textbf{merci} = \emph{mehr-SEE}, with the soft guttural French \emph{r} at the back of the throat (never the Spanish tap). Stress the end: \emph{mehr-\textbf{SEE}}.
\end{itemize}
\end{sounds}

\begin{cousinweb}
\textbf{merci} — \textquotedblleft{}thank you.\textquotedblright{} Root: Latin \textbf{mercēs}, \emph{\textquotedblleft{}wages, reward, price, that which is earned.\textquotedblright{}} Saying \emph{merci} originally invoked a \textbf{reward} owed for a kindness — and, in old religious use, a plea for divine \textbf{mercy} (a reward you haven't earned).

That one root \textbf{merc-} (\textquotedblleft{}pay, trade, reward\textquotedblright{}) pays out a huge English family:

\begin{itemize}
\raggedright
  \item \textbf{mercy}, \textbf{merciful} — the reward/pity sense.
  \item \textbf{mercenary} — a soldier who fights for \emph{wages}.
  \item \textbf{merchant}, \textbf{commerce}, \textbf{market}, \textbf{mercantile} — buying and selling, the \emph{price} sense (Latin \emph{mercārī}, \textquotedblleft{}to trade\textquotedblright{}).
  \item Even \textbf{Mercury} — the Roman god of merchants.
\end{itemize}
\end{cousinweb}

\begin{culture}
Here's a lovely contrast to lock in:

\noindent
\begin{tabularx}{\linewidth}{@{}>{\raggedright\arraybackslash}X>{\raggedright\arraybackslash}X>{\raggedright\arraybackslash}X@{}}
\toprule
\textbf{language} & \textbf{\textquotedblleft{}thank you\textquotedblright{}} & \textbf{literal root} \\
\midrule
\textbf{French} & \emph{merci} & \emph{mercēs} — \textbf{reward / mercy} \\
\textbf{Spanish} & \emph{gracias} & \emph{grātia} — \textbf{grace / favour} \\
\textbf{Portuguese} & \emph{obrigado} & \emph{obligātus} — \textquotedblleft{}I am \textbf{obliged}\textquotedblright{} \\
\bottomrule
\end{tabularx}

Three Romance sisters, three different pictures of gratitude: French hands you a \textbf{reward}, Spanish hands you \textbf{grace}, Portuguese declares itself \textbf{obliged}. Same feeling, three metaphors — this is exactly the kind of thing worth \emph{noticing}, not just memorizing.
\end{culture}

\subsection*{Your turn}
Say these aloud:

\begin{itemize}
\raggedright
  \item \textquotedblleft{}merci\textquotedblright{} — \emph{mehr-SEE}, guttural \emph{r}
  \item \textquotedblleft{}merci\textquotedblright{} then English \textquotedblleft{}mercy, merchant, commerce\textquotedblright{} — one family
  \item French \emph{merci} (reward) vs Spanish \emph{gracias} (grace) — feel the difference
\end{itemize}

\begin{tcolorbox}[breakable,colback=teal!4,colframe=teal!35!black,title={Before you move on}]
What everyday English word is \emph{merci} a cousin of? (\emph{Mercy}.) And its \textquotedblleft{}trade\textquotedblright{} cousins? (Merchant, commerce, market, mercenary.) How does its metaphor differ from Spanish \emph{gracias}? (\emph{merci} = reward/mercy; \emph{gracias} = grace.) Next: the reply — \textbf{de rien}.
\end{tcolorbox}
\section[de rien]{de rien — \textquotedblleft{}you're welcome,\textquotedblright{} or \textquotedblleft{}it was nothing\textquotedblright{}}

\label{lesson:FR-C03-de-rien}



\begin{quote}
Someone says \textbf{merci}. The easy reply: \textbf{de rien} — literally \emph{\textquotedblleft{}of nothing.\textquotedblright{}} You wave the favour off as no trouble — and Spanish pulls \emph{this exact trick} too.
\end{quote}

\begin{sounds}
\begin{itemize}
\raggedright
  \item \texttt{nasal-in} — \textbf{rien} is a nasal ending: the air goes out through the nose and there is no hard \emph{n} to stop it. \emph{duh RYAN}, with that final vowel humming rather than closing.
  \item \texttt{r-uvular} — that guttural French \emph{r} again, at the back of the throat.
\end{itemize}
\end{sounds}

\begin{cousinweb}
\begin{itemize}
\raggedright
  \item \textbf{de} = \textquotedblleft{}of / from\textquotedblright{} (Latin \textbf{dē}).
  \item \textbf{rien} = \textquotedblleft{}nothing.\textquotedblright{} And here is the beautiful part.
\end{itemize}

\textbf{rien} comes from Latin \textbf{rem} — the accusative of \textbf{rēs}, \emph{\textquotedblleft{}a thing.\textquotedblright{}} In Latin \emph{rem} meant a \textbf{thing}, something real. But French kept using it inside negations — \emph{ne … rien}, \textquotedblleft{}not … a thing\textquotedblright{} — until the \emph{thing} word absorbed the \textquotedblleft{}nothing\textquotedblright{} and \textbf{rien came to mean 'nothing' all by itself.}

Latin \emph{rēs} also gives English \textbf{real}, \textbf{re}public (\emph{rēs pūblica}, \textquotedblleft{}the public thing\textquotedblright{}), and \textbf{rebus}. So \emph{rien} is, at the root, a \textbf{real thing} that learned to mean nothing at all.
\end{cousinweb}

\begin{culture}
If that sounds familiar, it should: \textbf{Spanish did the identical thing with nada} ($\leftarrow$ Latin \emph{nāta}, \textquotedblleft{}a born thing\textquotedblright{} $\to$ \textquotedblleft{}nothing\textquotedblright{}). Two sister languages, the same move — the word for \textbf{thing} collapses into the word for \textbf{nothing}:

\noindent
\begin{tabularx}{\linewidth}{@{}>{\raggedright\arraybackslash}X>{\raggedright\arraybackslash}X>{\raggedright\arraybackslash}X@{}}
\toprule
\textbf{language} & \textbf{\textquotedblleft{}you're welcome\textquotedblright{}} & \textbf{the \textquotedblleft{}nothing\textquotedblright{} word, from Latin} \\
\midrule
\textbf{French} & \emph{de rien} & \emph{rien} $\leftarrow$ \emph{rem}, \textquotedblleft{}a thing\textquotedblright{} \\
\textbf{Spanish} & \emph{de nada} & \emph{nada} $\leftarrow$ \emph{nāta}, \textquotedblleft{}a born thing\textquotedblright{} \\
\bottomrule
\end{tabularx}

Two sister languages, one instinct: wave the favour off as no \emph{thing} at all.
\end{culture}

\subsection*{Your turn}
Say these aloud:

\begin{itemize}
\raggedright
  \item \textquotedblleft{}merci\textquotedblright{} … \textquotedblleft{}de rien\textquotedblright{} — the full exchange, both voices
  \item French \emph{de rien} and Spanish \emph{de nada} back to back — same idea, two \textquotedblleft{}thing $\to$ nothing\textquotedblright{} words
\end{itemize}

\begin{tcolorbox}[breakable,colback=teal!4,colframe=teal!35!black,title={Before you move on}]
What did \emph{rien} (\textquotedblleft{}nothing\textquotedblright{}) originally mean in Latin? (\emph{rem} — a \emph{thing}.) What Spanish word pulled the same trick, and from what? (\emph{nada} $\leftarrow$ \emph{nāta}, \textquotedblleft{}a born thing.\textquotedblright{}) Next: the verb French uses for \textquotedblleft{}how are you?\textquotedblright{} — and it's not \textquotedblleft{}to stand\textquotedblright{} like Spanish, but \textbf{to go}.
\end{tcolorbox}
\section[aller]{aller — \textquotedblleft{}to go,\textquotedblright{} the verb behind \textquotedblleft{}how are you?\textquotedblright{}}

\label{lesson:FR-C03-aller}



\begin{quote}
Spanish asks \emph{how are you} with \textbf{estar}, \textquotedblleft{}to \emph{stand}.\textquotedblright{} French asks it with \textbf{aller}, \textquotedblleft{}to \textbf{go}\textquotedblright{} — \emph{how does it go?} Same question, a different picture of a life: one language stands, the other walks. Meet \textbf{aller}.
\end{quote}

\begin{sounds}
\begin{itemize}
\raggedright
  \item \texttt{er-ending} — the infinitive \textbf{-er} is a clean \emph{ay}: \textbf{aller} = \emph{ah-LAY} (the double \emph{l} is just one \emph{l} sound; the \emph{r} is silent here).
\end{itemize}
\end{sounds}

\begin{cousinweb}
\textbf{aller} means \textquotedblleft{}to go.\textquotedblright{} Its origin is one of the great messes of French — it is \textbf{suppletive}, meaning it was stitched together from \textbf{three} different Latin verbs, and you can still see the seams in its everyday forms:

\begin{itemize}
\raggedright
  \item \textbf{je vais} (\textquotedblleft{}I go\textquotedblright{}) $\leftarrow$ Latin \textbf{vādere}, \textquotedblleft{}to go, to rush\textquotedblright{} $\to$ English \textbf{invade}, \textbf{evade}, \textbf{wade} (all \textquotedblleft{}going\textquotedblright{} words).
  \item \textbf{nous allons} (\textquotedblleft{}we go\textquotedblright{}), and the infinitive \textbf{aller} $\leftarrow$ Latin \textbf{ambulāre}, \textquotedblleft{}to walk\textquotedblright{} $\to$ English \textbf{amble}, \textbf{ambulance} (a walking hospital), \textbf{preamble}, \textbf{perambulate}.
  \item \textbf{j'irai} (\textquotedblleft{}I will go\textquotedblright{}) $\leftarrow$ Latin \textbf{īre}, \textquotedblleft{}to go\textquotedblright{} $\to$ English \textbf{exit} (\emph{ex-it}, \textquotedblleft{}he goes out\textquotedblright{}), \textbf{transit}, \textbf{initial}.
\end{itemize}

So a single French verb carries \textbf{three} Latin roots for motion. You don't need to conjugate it yet — just recognize one form for the next lesson:

\begin{itemize}
\raggedright
  \item \textbf{ça va} = \textquotedblleft{}it goes\textquotedblright{} (the \emph{va} is from \emph{vādere}) — the beating heart of French small talk.
\end{itemize}
\end{cousinweb}

\begin{grammarlens}[title={Grammar Lens: \textquotedblleft{}going\textquotedblright{} as \textquotedblleft{}being\textquotedblright{}}]
Lots of languages borrow a motion verb to talk about \emph{state}: English says \textquotedblleft{}how are you \textbf{getting on}? / how's it \textbf{going}?\textquotedblright{} French made that the standard: \emph{Comment ça va?} = \textquotedblleft{}How does it \textbf{go}?\textquotedblright{} Hold the metaphor — it makes the next lesson obvious.
\end{grammarlens}

\subsection*{Your turn}
Say these aloud:

\begin{itemize}
\raggedright
  \item \textquotedblleft{}aller\textquotedblright{} — \emph{ah-LAY}
  \item \textquotedblleft{}ça va\textquotedblright{} — \emph{sa VA}, \textquotedblleft{}it goes\textquotedblright{}
  \item English \textquotedblleft{}amble, ambulance\textquotedblright{} (walk) and \textquotedblleft{}invade, evade\textquotedblright{} (go) — aller's two big families
\end{itemize}

\begin{tcolorbox}[breakable,colback=teal!4,colframe=teal!35!black,title={Before you move on}]
What does \emph{aller} literally mean, and how many Latin verbs is it stitched from? (To \emph{go}; three — \emph{ambulāre, vādere, īre}.) Which everyday form means \textquotedblleft{}it goes\textquotedblright{}? (\emph{ça va}.) So \emph{Comment ça va?} literally asks — what? (\textquotedblleft{}How does it \emph{go}?\textquotedblright{}) Next: ask it for real.
\end{tcolorbox}
\section[comment ça va ?]{Comment ça va ? — “how are you?”}

\label{lesson:FR-C03-comment-ca-va}



\begin{quote}
You know \textbf{comment}, “how,” and \textbf{ça va}, “it goes.” Assemble the most useful neutral check-in in spoken French.

\begin{quote}
\textbf{Comment ça va ?} — “How are you?” literally: “How does this go?”
\end{quote}

The even shorter \textbf{Ça va ?} is extremely common. Answer with the same words:

\begin{quote}
\textbf{Ça va ? — Ça va.} “How’s it going? — It’s going / I’m fine.”
\end{quote}
\end{quote}

\begin{cousinweb}
\begin{itemize}
\raggedright
  \item \textbf{comment} comes from Latin \emph{quōmodo}, “in what manner”;
  \item \textbf{ça} comes from \emph{ecce hoc}, “behold this”;
  \item \textbf{va} is the third-person form of \textbf{aller}, “to go.”
\end{itemize}

French pictures wellbeing as \textbf{going}. Spanish uses a “standing” verb instead; the comparison is a memory hook, not a prerequisite.

Pronounce \emph{comment} roughly \emph{ko-mahn}: the ending is nasal and the written \emph{t} is silent in this neutral phrase.
\end{cousinweb}

\subsection*{Your turn}
Say these aloud:

\begin{itemize}
\raggedright
  \item “Comment ça va ?”
  \item the shorter question — “Ça va ?”
  \item the pair — “Ça va ? — Ça va.”
\end{itemize}

\emph{Twice through:} “Ça va ? — Ça va.”

\begin{tcolorbox}[breakable,colback=teal!4,colframe=teal!35!black,title={Before you move on}]
What does \emph{Comment ça va ?} literally ask? (“How does it go?”) Which verb owns \textbf{va}? (\textbf{Aller}.) What is the shortest version? (\textbf{Ça va ?}) How can you answer it? (\textbf{Ça va.}) Next: choose \emph{tu} or \emph{vous} explicitly.
\end{tcolorbox}
\section[Comment vas-tu ? / Comment allez-vous ?]{Comment vas-tu ? / Comment allez-vous ?}

\label{lesson:FR-C03-comment-registers}



\begin{quote}
\textbf{Comment ça va ?} avoids naming the listener. Now use the \emph{tu / vous} choice you learned in Chapter 2.
\end{quote}

\begin{grammarlens}[title={Grammar Lens: a ladder of three registers}]
French offers a neat climb from casual to formal. \emph{tu} is the familiar “you” (a friend, a peer); \emph{vous} is the formal or plural “you” (a stranger, an elder, a group).

\noindent
\begin{tabularx}{\linewidth}{@{}>{\raggedright\arraybackslash}X>{\raggedright\arraybackslash}X>{\raggedright\arraybackslash}X@{}}
\toprule
\textbf{register} & \textbf{question} & \textbf{literally} \\
\midrule
casual, universal & \textbf{Comment ça va ?} (or \textbf{Ça va ?}) & “How does \emph{it} go?” \\
informal, one friend & \textbf{Comment vas-tu ?} & “How go \emph{you}?” \\
formal or plural & \textbf{Comment allez-vous ?} & “How go \emph{you} (vous)?” \\
\bottomrule
\end{tabularx}

\textbf{Vas} and \textbf{allez} are simply the \emph{tu} and \emph{vous} forms of \emph{aller}, “to go.” Formal writing and careful speech invert verb and pronoun: \textbf{allez-vous}, literally “go-you.” For a stranger or a boss, reach for \textbf{Comment allez-vous ?}
\end{grammarlens}

\begin{sounds}
The final \emph{t} of \textbf{comment} is normally silent. Before vowel-initial \textbf{allez}, liaison can carry it across:

\begin{quote}
\textbf{Comment allez-vous ?} — roughly \emph{ko-mahn-ta-lay-voo}
\end{quote}

That linking sound is not an extra word; it is a normally silent consonant surfacing between two words.

Use \textbf{Ça va ?} as the easy neutral default, \textbf{vas-tu} for a friend when you want explicit informality, and \textbf{allez-vous} for formal or plural address.
\end{sounds}

\subsection*{Your turn}
Say these aloud:

\begin{itemize}
\raggedright
  \item to a friend — “Comment vas-tu ?”
  \item to a stranger — “Comment allez-vous ?”
  \item the neutral default — “Ça va ?”
\end{itemize}

\emph{Twice through:} “Vas-tu — friend. Allez-vous — formal or plural.”

\begin{tcolorbox}[breakable,colback=teal!4,colframe=teal!35!black,title={Before you move on}]
Which form addresses one friend? (\textbf{Comment vas-tu ?}) Which is formal or plural? (\textbf{Comment allez-vous ?}) What sound appears in the liaison? (The normally silent \emph{\textbf{t}} of \emph{comment}.) Next: the “so-so” answer.
\end{tcolorbox}
\section[comme ci, comme ça]{comme ci, comme ça — the classic French shrug}

\label{lesson:FR-C03-comme-ci-comme-ca}



\begin{quote}
Someone asks \emph{Ça va?} You're not great, not bad. The perfect, gloriously French answer — a phrase English borrowed whole — is \textbf{comme ci, comme ça}: \textquotedblleft{}like this, like that.\textquotedblright{} A verbal shrug of the shoulders.

First, twenty seconds of review. From Chapter 1: \textbf{bien} — \textquotedblleft{}well\textquotedblright{} ($\leftarrow$ Latin \emph{bene}). So \emph{Ça va bien} = \textquotedblleft{}It's going well.\textquotedblright{} Now the answer for when it's only \emph{going}.
\end{quote}

\begin{cousinweb}
\textbf{comme ci, comme ça} — literally \textbf{\textquotedblleft{}like this, like that.\textquotedblright{}} Every piece is an old friend:

\begin{itemize}
\raggedright
  \item \textbf{comme} = \textquotedblleft{}like, as\textquotedblright{} — from Latin \textbf{quōmodo}, \textquotedblleft{}in what manner.\textquotedblright{} That's the \textbf{same root as \emph{comment}} (\textquotedblleft{}how\textquotedblright{}) from last lesson! \emph{comment} = \textquotedblleft{}in what manner (as a question)\textquotedblright{}; \emph{comme} = \textquotedblleft{}in the manner of, like.\textquotedblright{} One Latin word, two French children.
  \item \textbf{ci} = \textquotedblleft{}here / this\textquotedblright{} $\leftarrow$ Latin \textbf{ecce hic}, \textquotedblleft{}behold here.\textquotedblright{}
  \item \textbf{ça} = \textquotedblleft{}there / that\textquotedblright{} $\leftarrow$ Latin \textbf{ecce hoc}, \textquotedblleft{}behold this\textquotedblright{} — the same \emph{ça} from \emph{Ça va?}.
\end{itemize}

So you're literally waving a hand: \emph{a bit like this way, a bit like that way} — neither up nor down.
\end{cousinweb}

\begin{grammarlens}[title={Grammar Lens: your answer set for \emph{Ça va?}}]
\noindent
\begin{tabularx}{\linewidth}{@{}>{\raggedright\arraybackslash}X>{\raggedright\arraybackslash}X@{}}
\toprule
\textbf{answer} & \textbf{sense} \\
\midrule
\textbf{(Très) bien, merci} & (very) well, thank you \\
\textbf{Ça va} & fine / okay (\textquotedblleft{}it goes\textquotedblright{}) \\
\textbf{Comme ci, comme ça} & so-so \\
\textbf{Pas mal} & not bad \\
\bottomrule
\end{tabularx}
\end{grammarlens}

\begin{culture}
Every language has its shrug, and each paints it differently: Spanish says \textbf{más o menos} (\textquotedblleft{}more or less\textquotedblright{}) or \textbf{regular}, Italian says \textbf{così così} (\textquotedblleft{}so so\textquotedblright{}), French says \textbf{comme ci, comme ça} (\textquotedblleft{}like this, like that\textquotedblright{}). Same gesture, three pictures.
\end{culture}

\subsection*{Your turn}
Say these aloud:

\begin{itemize}
\raggedright
  \item \textquotedblleft{}comme ci, comme ça\textquotedblright{} — with a little hand-wobble in the voice
  \item answer \textquotedblleft{}Ça va?\textquotedblright{} three ways — \textquotedblleft{}Bien\textquotedblright{} / \textquotedblleft{}Ça va\textquotedblright{} / \textquotedblleft{}Comme ci, comme ça\textquotedblright{}
  \item notice \emph{comme} and \emph{comment} share a root — \emph{quōmodo}
\end{itemize}

\begin{tcolorbox}[breakable,colback=teal!4,colframe=teal!35!black,title={Before you move on}]
What does \emph{comme ci, comme ça} literally mean? (\textquotedblleft{}Like this, like that.\textquotedblright{}) What \textquotedblleft{}how\textquotedblright{}-word shares \emph{comme}'s root? (\emph{comment} — both from \emph{quōmodo}.) What's the Spanish equivalent shrug? (\emph{más o menos} / \emph{regular}.) Next: put the whole exchange together.
\end{tcolorbox}
\section[Practice]{Practice — \textquotedblleft{}comment ça va ?\textquotedblright{}, start to finish}

\label{lesson:FR-C03-practice}



\begin{quote}
No new words. This chapter's atoms — \emph{merci, de rien, aller, ça va, comme ci comme ça} — assemble into a real exchange, run formally and informally. It picks up right where the Chapter 2 introduction left off.
\end{quote}

\subsection*{The exchange}
\textbf{Formal} (a stranger, an elder):

\begin{quote}
— Bonjour. \textbf{Comment allez-vous ?} — \textbf{Très bien, merci.} Et vous ? — \textbf{Comme ci, comme ça.} — {[}later, after a small favour{]} — Merci ! — \textbf{De rien.}
\end{quote}

\textbf{Informal} (a friend):

\begin{quote}
— Salut ! \textbf{Ça va ?} — \textbf{Ça va, et toi ?} — Comme ci, comme ça.
\end{quote}

The glue word is \textbf{et vous ? / et toi ?} — \textquotedblleft{}and you?\textquotedblright{} (\emph{et} = \textquotedblleft{}and,\textquotedblright{} Latin \emph{et}; \emph{toi} is the stressed \textquotedblleft{}you\textquotedblright{}). It bounces the question back — exactly like Spanish \emph{¿y usted? / ¿y tú?}

\begin{grammarlens}[title={What you've built this chapter}]
\begin{itemize}
\raggedright
  \item \textbf{merci} — thank you ($\leftarrow$ \emph{mercēs}, \textquotedblleft{}reward / mercy\textquotedblright{}).
  \item \textbf{de rien} — you're welcome ($\leftarrow$ \emph{rem}, \textquotedblleft{}a thing\textquotedblright{} $\to$ \textquotedblleft{}nothing\textquotedblright{}; twin of Spanish \emph{de nada}).
  \item \textbf{aller / ça va} — \textquotedblleft{}to go,\textquotedblright{} the verb French uses for \emph{how are you} (vs Spanish \emph{estar}, \textquotedblleft{}to stand\textquotedblright{}).
  \item \textbf{Comment ça va ? / allez-vous ? / vas-tu ?} — the casual, formal, and informal question.
  \item \textbf{bien / ça va / comme ci comme ça / pas mal} — the answer ladder.
\end{itemize}
\end{grammarlens}

\subsection*{Your turn}
Say these aloud:

\begin{itemize}
\raggedright
  \item the full \emph{formal} exchange, both voices
  \item the same \emph{informally} — \emph{allez-vous $\to$ ça va}, \emph{vous $\to$ toi}
  \item chain Chapter 2 + 3 — \textquotedblleft{}Je m'appelle … Comment allez-vous ?\textquotedblright{}
\end{itemize}

\emph{Twice through:} \textquotedblleft{}Ça va ? — Ça va, et toi ?\textquotedblright{} until it's reflex.

\begin{tcolorbox}[breakable,colback=teal!4,colframe=teal!35!black,title={Before you move on}]
Which verb carries French \textquotedblleft{}how are you?\textquotedblright{}, and how does its metaphor differ from Spanish's? (\emph{aller}, \textquotedblleft{}to go\textquotedblright{} — \emph{how does it go?} — vs \emph{estar}, \textquotedblleft{}to stand.\textquotedblright{}) What does \emph{de rien} literally say, and its Spanish twin? (\textquotedblleft{}Of nothing\textquotedblright{}; \emph{de nada}.) Next chapter: \textbf{farewells} — au revoir, à bientôt, bonne journée.
\end{tcolorbox}
